
d'où le résultat d'%
après (6.1.2) en passant au quotient.

Si \(a\in K\), \(\widetilde c(a)=(a,1)\), donc \(\widetilde c\) est un homomorphisme de \(K\)-algèbres~; de plus il est compatible aux augmentations.

Il reste enfin à voir que c'est un \(\lambda\)-homomorphisme. Il faut donc montrer la commutativité du diagramme :
\[
\begin{tikzcd}[column sep=large,row sep=large]
A \arrow[r,"\gamma_t"] \arrow[d,"\widetilde c"'] &
\widehat G_g(A) \arrow[d,"\widehat G_g(\widetilde c)"]\\
\Ch(\Gr A) \arrow[r,"\gamma_t"] &
\widehat G_g(\Ch(\Gr A))
\end{tikzcd}
\]

Par additivité, il suffit de la montrer pour les éléments de \(K\), pour lesquels elle est évidente, et pour les éléments d'augmentation nulle de \(A\), pour lesquels elle résulte, d'après la définition de \(c\) et (6.1.4), de celle du diagramme (3.7.2).

\textbf{Proposition 6.9.} \uline{Soit} \(A\) \uline{une} \(K\)-\(\lambda\)-\uline{algèbre augmentée}, \(z\in\Fil^mA\)~; \uline{alors} :
\[
  \gamma^m(z)\equiv(-1)^{m-1}(m-1)!z\ \mathrm{mod.}\ \Fil^{m+1}A,
\]
\uline{i}.\uline{e}.
\[
  c^m(z)=\cl_{\Gr^mA}\!\left((-1)^{m-1}(m-1)!z\right).
\]

Il faut voir que ces deux éléments donnent même image dans \(\Gr^mA\), c'est-à-dire que la \(m\)-ième classe de Chern de \(z\) est l'image de \((-1)^{m-1}(m-1)!z\).

On peut écrire \(z\) comme une somme \(\sum a\mathbin{.}\gamma^{i_1}(x_1)\ldots\gamma^{i_k}(x_k)\), avec \(a\in A\), \(\epsilon(x_i)=0\) pour tout \(i\), et \(i_1+\ldots+i_k\ge m\). On a alors :
\[
  \widetilde c(z)=
  \sum \widetilde c(a)\mathbin{.}
        \widetilde c(\gamma^{i_1}(x_1))\ldots
        \widetilde c(\gamma^{i_k}(x_k))
\]
Comme \(\widetilde c\) est un \(\lambda\)-homomorphisme, on peut écrire :
\[
  \widetilde c(\gamma^{i_j}(x_j))=\gamma^{i_j}(\widetilde c(x_j)).
\]
